\documentclass[11pt]{article}
% Diplomatic transcription macros (12 Nov 2026). Typed text is plain; the handwritten layer is blue,
% deletions red, typewriter strike-outs grey, doubtful readings in brackets with a question mark.
\usepackage[utf8]{inputenc}
\usepackage[T1]{fontenc}
\usepackage{lmodern}
\usepackage[a4paper,margin=2.2cm]{geometry}
\usepackage{amsmath,amssymb}
\usepackage{xcolor}
\usepackage{tikz-cd}
\usepackage{environ}
\setlength{\parindent}{0pt}
\setlength{\parskip}{4pt}
\title{EGA V, English translation by P. Blass, 1987 (cote 24)\\[4pt]\large Diplomatic transcription from the Montpellier scans}
\author{AI-transcribed by Claude Opus 5.5}
\date{12 November 2026}
\newcommand{\scanpage}[1]{\par\bigskip\noindent{\small\sffamily[scan #1]}\par}
\newcommand{\mspage}[1]{{\small\sffamily[p.\,#1]}}
% The markup may stand in text or in math, and its argument may itself contain $...$: in math mode an argument
% with a $ in it is set through \text (so that its $...$ is math again), one without stays math.
\makeatletter
\long\def\dip@has@dollar#1$#2\dip@end{\ifx\relax#2\relax\@secondoftwo\else\@firstoftwo\fi}
\newcommand{\dip@wrap}[1]{\ifmmode\dip@has@dollar#1$\dip@end{\text{#1}}{#1}\else#1\fi}
\newcommand{\hand}[1]{\textcolor{blue}{\dip@wrap{#1}}}
\let\over\relax
\newcommand{\over}[2]{\textcolor{gray}{\dip@wrap{#1}}\textcolor{blue}{\ensuremath{^{\!\nearrow}}\dip@wrap{#2}}}
\newcommand{\inserted}[2][]{\textcolor{blue}{\ensuremath{\langle}\dip@wrap{#2}\ensuremath{\rangle}}}
\newcommand{\margin}[1]{\textcolor{blue}{\{\dip@wrap{\textsf{marg.:} #1}\}}}
\newcommand{\deleted}[1]{\textcolor{red}{[\dip@wrap{\textsf{del.:} #1}]}}
\newcommand{\xout}[1]{\textcolor{gray}{[\dip@wrap{\textsf{x:} #1}]}}
\newcommand{\uline}[1]{\underline{\dip@wrap{#1}}}
\makeatother
\newcommand{\uuline}[1]{\underline{\underline{#1}}}
\newcommand{\unsure}[1]{[#1?]}
\newcommand{\illegible}[1]{[\textsf{illeg.\ #1}]}
\newcommand{\diagram}[1]{\par{\centering #1\par}}
\newenvironment{diagramblock}{\par\begin{center}}{\end{center}\par}

\begin{document}
\maketitle

% EGA V, English translation by P. Blass, 1987 (cote 24): archive cote 24, site page 4 (PDF page 5).
% typescript. Draft compared: none.
% AI-transcribed by Claude Opus 5.5 (claude-opus-5-5, adaptive thinking, effort high), Message Batches API, batch pilot (msgbatch_0189Nhf1xB615YWnztubtLfc), 2026-11-12; not yet checked by eye.
\scanpage{24/4}
\mspage{4}

\uline{HYPERPLANE SECTIONS AND CONIC PROJECTIONS}

\uline{Section 1: Preliminaries and Notation.}

Let S be a prescheme, let $\underline{E}$ be a locally free module of
finite type over S let $\underline{E}^{\vee}$ be its dual. We denote by P = $\underline{P}(\underline{E})$ the
projective fibration defined by $\underline{E}$ and by P$^{\vee}$ the projective
fibration defined by E$^{\vee}$. We shall call P$^{\vee}$ the \uline{scheme of}
\uline{hyperplanes} of P. This terminology can be justified as follows.
Let $\xi$ be a section of P$^{\vee}$ over S which is therefore determined by an
invertible quotient module L of $\underline{E}^{\vee}$. From it we obtain an
invertible quotient L$_{P}$ of $(\underline{E}^{\vee})_{\underline{P}} = (\underline{E}_{P})^{\vee}$, on the other hand, we have
the invertible quotient module O$_{P}$(1) of E$_{P}$. Passing to duals we
may take $\underline{L}_{P}{-1}$ resp. $\underline{O}_{P}(-1)$ to be invertible submodules (locally
direct \uline{summands or factors}) of $\underline{E}_{P}$ (resp. of $(\underline{E}_{P})^{\vee}$) and the pairing
$EP \otimes E^{\vee}_{P} \to \underline{O}_{P}$ defines therefore a natural pairing (*) $\underline{O}_{P}(-1) \otimes \underline{L}_{P}{-1}$
$\to \underline{O}_{P}$ or also the transpose homomorphism (**) $\underline{O}_{P} \to \underline{O}_{P}(1) \otimes \underline{L}_{P} =$
$\underline{O}_{P}(1)$ i.e. a section of $\underline{L}_{P}(1)$ canonically defined by $\xi$. The
``divisor'' of that section i.e. the closed subscheme H$_{\xi}$ of $\underline{P}$ defined
by the image ideal of (*) is called the hyperplane in $\underline{P}$ defined by
the element $\xi \in$ P$^{\vee}$(S). We could describe it by noting that locally
over $\underline{S}$, $\xi$ is given by a section $\phi$ of E such that $\phi$ (s) $\neq 0$ for all
s ($\phi$ is determined by $\xi$ up to multiplication by an invertible
section of $\underline{O}_{S}$); since E = p$^{*}$(O$_{P}$(1)) (p: P $\to$ S being the
projection), $\phi$ can be considered as a section of O$_{P}$(1) the divisor
of which is nothing else but H$_{\xi}$.

Of course, if we consider L-1 as an invertible submodule of E
locally a direct factor in E then the correspondence between $\xi$
% NOTE: a white strip covers the next typed line (only faint traces show through); nothing else is typed below it on the page

\clearpage
% EGA V, English translation by P. Blass, 1987 (cote 24): archive cote 24, site page 18 (PDF page 19).
% typescript. Draft compared: none.
% AI-transcribed by Claude Opus 5.5 (claude-opus-5-5, adaptive thinking, effort high), Message Batches API, batch pilot (msgbatch_0189Nhf1xB615YWnztubtLfc), 2026-11-12; not yet checked by eye.
\scanpage{24/18}
\mspage{18}

not be immersed it is necessary and sufficient that $Z_i$ should be
such [PB incomprehensible].* If that is so then the radicial
multiplicity of $F$ at $Z_i$ at $k$ is equal to that of $G_\eta$ at $Z_i$ relative
to $k(\eta)$.

The first assertion is contained in the last assertion of
2.7. For the second, since $Y_\eta \rightarrow X$ is flat therefore the functor $F$
$\rightarrow G_\eta$ is exact, we are reduced by restriction to a neighborhood of
the general point of $Z_i$ and by a devissage (unscrewing) to the case
where $F = O_{Z_i}$ and we may assume $Z_i = X$. Also we could start by
assuming that $X$ is separable over $k$ is reduced to the case of $k$
algebraically closed [PB: incomprehensible].* Then the assertion
is contained in 2.15. (iii). Then we conclude as usual:

\uline{Corollary 2.17}

Let $Z$ be the set of points of $X$ where $F$ is not separable over
$k$ (reference). Then $G_\eta$ is separable over $k(\eta)$ if and only if $Z$ is
finite. In particular if $F$ is separable over $k$ then $G_\eta$ is
separable over $k(\eta)$.

\uline{Remark 2.18}

The key result 2.12 (and its consequences 2.13 and 2.17)
become false if we abandon the assumption that $f$ is an immersion
as we see for example by taking for $X$ a purely inseparable
covering of $\underline{P}$. However if $k$ is of characteristic zero 2.12 and
2.17 are valid without assuming that $f$ is an immersion.

\clearpage
% EGA V, English translation by P. Blass, 1987 (cote 24): archive cote 24, site page 37 (PDF page 38).
% typescript. Draft compared: none.
% AI-transcribed by Claude Opus 5.5 (claude-opus-5-5, adaptive thinking, effort high), Message Batches API, batch pilot (msgbatch_0189Nhf1xB615YWnztubtLfc), 2026-11-12; not yet checked by eye.
\scanpage{24/37}
\mspage{37}

have dim f($X_{i}$) > [illegible] [is it two ask Grothendieck]* then Z
is of codimension two in $\underline{P}$v.  This follows from 2.1 and 2.2.
(which implies that every irreducible component of $\underline{Y}$ dominates $\underline{P}$v)
and from the dimension theory for the morphism $\underline{Y} \rightarrow \underline{P}$v.  Starting
from this result we may give as a corollary the case where we
start from a closed subset Z of X and where we consider the
dimension of the inverse images $Z_{\xi}$ in the $Y_{\xi}$ ($\xi \in \underline{P}$v) and we may
even take for Z the set of $\xi \in \underline{P}$v such that there exists an
irreducible component of $T_{k(\xi)}$ whose trace on $Y_{\xi}$ has the greatest
dimension (NB we always assume that for every irreducible $Z_{i}$ of Z
we have dim f($Z_{i}$) > 0).

Finally the most precise statement in this direction and one
that results easily from the first announcement (for X irreducible) and
from 2.7 is the following modified statement:  F being coherent
over X, suppose that for every associated prime cycle T for F we
have dim f(T) > 0 then the set of $\xi \in \underline{P}$v such that $\phi_{\xi}$ is not $F_{k(\xi)}$ -
regular is (constructible and) of codimension $\geq$ 2.  (The notation
for $\phi_{\xi}$ is that from N05).  We can give this as the principal
assertion, and announce the previous assertions as corollaries,
the proof being or proceeding via one of the corollaries.

Please note that with the preceding notations if $\xi \in \overset{\mathrm{v}}{\underline{P}}$ - Z,
then for every y $\in \underline{Y}_{\xi}$ we have coprof$_{y}$ $G_{k(\xi)}$ = coprof$_{y}$ $G_{\xi}$ and
consequently if coprof F $\leq$ n in particular if F is Cohen-Macauley
then for $\xi \in \underline{P}$v - Z, $G_{\xi}$ is Cohen-Macauley.  Finally if F is ($S_{k}$)
[illegible ask Grothendieck]*

$S_{k-1}$ for $\xi \in \underline{P}$v - Z (reference $O_{IV}$).

\margin{* Ask AG.}

\end{document}
